% Section 5.5, from lectures/L05/appendix/d_exercises.md.

\section{Exercises}
\label{c5:sec:exercises}

\begin{aside}
\textbf{How to check your work.} Exercise~\ref{c5:ex:driver} is checked by this chapter's test
suite: write the file at the path the specification gives, then run \code{make test}. The suite is
cumulative and runs Chapter~1 to Chapter~4's tests too.

The rest will have worked solutions in the companion repository, published later as
\file{lectures/L05/appendix/e_solutions.md}, and the hand calculations have short answers in
Appendix~\ref{app:answers}. Work each one before you read it.

\textbf{No hardware is needed for any of this.}
\end{aside}

Each exercise is labelled with its kind. There is exactly one \textbf{Cross-check}, and here it is
Exercise~\ref{c5:ex:crosscheck}.

% ------------------------------------------------------------------------------------------
\exercise{Why not a loop}{Recall}
\label{c5:ex:loop}
\begin{parts}
  \item Give three reasons a counting loop is a worse way to wait than a timer. One of them is
    about the processor, one about arithmetic, and one about interrupts.
  \item Name the four pieces of hardware in the timer and say which two you choose.
  \item What does CTC stand for, and what does the hardware do on a match that normal mode does
    not?
  \item A delay loop and a CTC timer are each one cycle out per period. After an hour, how do the
    two errors compare? Say why.
\end{parts}

% ------------------------------------------------------------------------------------------
\exercise{Prescaler and compare}{Hand calculation}
\label{c5:ex:compare}
At 16\,MHz, for each frequency below, give the prescaler and the compare value for a 16-bit timer
and for an 8-bit one, and the error in each case.
\begin{parts}
  \item 500\,Hz
  \item 440\,Hz
  \item 3\,kHz
  \item For one of the three, both timers are exact. Say which and why it is not luck.
  \item For another, the 8-bit timer is more than an order of magnitude worse. Say which, give the
    ratio, and explain it in terms of ticks rather than of bits.
\end{parts}

\noindent\textbf{Check yourself:} the 3\,kHz row is worked in \secref{c5:sec:actual}, and the rule
that produces every row is \secref{c5:sec:choosing}.

% ------------------------------------------------------------------------------------------
\exercise{Choosing a prescaler}{Design}
\label{c5:ex:choose}
\begin{parts}
  \item State the rule for choosing a prescaler, and justify it in one sentence.
  \item The ratios are 1, 8, 64, 256, 1024. A frequency needs 70000 ticks at a prescaler of 1.
    Which prescaler do you end up at, how many ticks is that, and how much resolution did you lose
    compared with the one that did not fit?
  \item The four steps in that ladder are not all the same size. Say which are which, and when you
    would notice.
  \item Somebody suggests always using a prescaler of 1024 ``to be safe''. Give two things that goes
    wrong.
\end{parts}

% ------------------------------------------------------------------------------------------
\exercise{The registers}{Recall}
\label{c5:ex:registers}
\begin{parts}
  \item Name the four registers Timer1 needs for a CTC interrupt, and what goes in each.
  \item All four are at \code{0x60} or above. What does that rule out, and is the failure loud or
    quiet?
  \item \code{TCCR1B} holds both the mode bit and the clock select bits. Describe the bug that
    follows from that, and what it looks like from outside.
  \item The datasheet says to write a 16-bit register's high byte first. Say why, and say whether
    the simulator would catch you getting it wrong.
\end{parts}

% ------------------------------------------------------------------------------------------
\exercise{The software timer}{Code}
\label{c5:ex:driver}
Write \file{drivers/source/timer.asm} to the specification in \secref{c5:sec:driver}.
\begin{parts}
  \item Write all five subroutines.
  \item Run \code{make test}.
  \item Change \code{timer_tick} so that it compares before incrementing rather than after. Note
    which tests fail, then put it back.
  \item Change it to compare only the low bytes of the count and the target. Note which single test
    catches it, and say why every other test passes.
\end{parts}

% ------------------------------------------------------------------------------------------
\exercise{The hardware setup}{Code}
\label{c5:ex:setup}
Extend \file{drivers/app/main.asm} to the specification in \secref{c5:sec:setup}.
\begin{parts}
  \item Configure Timer1 for a 1\,kHz compare-match interrupt, using the prescaler and compare
    value you worked out by hand.
  \item Add the handler at the \code{TIMER1_COMPA} vector. Give the \code{.org} address you used
    and say where the number came from.
  \item Initialise a second LED, on Arduino pin 8, and have the handler call \code{timer_tick} on a
    software timer with a target that blinks it at a rate you can see. Say what target you chose
    and why.
  \item \code{make build} should build \file{app.hex}.
\end{parts}

% ------------------------------------------------------------------------------------------
\exercise{What cannot be reached}{Hand calculation}
\label{c5:ex:reach}
\begin{parts}
  \item Give the formula for the slowest frequency a timer can reach, and evaluate it at 16\,MHz for
    both counter widths.
  \item Can a 16-bit timer produce 0.5\,Hz? Give the prescaler and compare value, or say why not.
  \item Can it produce 0.1\,Hz? Same question.
  \item An 8-bit timer cannot produce 1\,Hz. Explain in one sentence why no value of \code{OCR0A}
    helps.
  \item You need 0.1\,Hz and you have a 16-bit timer. Describe the arrangement you would use, and
    give one interrupt rate and one target that achieve it.
  \item What is the fastest compare rate any of these timers can produce, and why is asking for
    more not merely inaccurate?
\end{parts}

% ------------------------------------------------------------------------------------------
\exercise[fillaccent]{How far apart are two interrupts}{Cross-check}
\label{c5:ex:crosscheck}
\emph{Compute it by hand, measure it, reconcile.}
\begin{parts}
  \item \textbf{By hand.} For a 1\,kHz compare-match interrupt from Timer1 at 16\,MHz, give the
    prescaler, the compare value, and the number of CPU cycles between two interrupts.
  \item \textbf{In cycles.} Turn your \code{OCR1A} back into a period:
    $\text{prescaler} \times (\mathrm{OCR1A} + 1)$ cycles, and then into microseconds. Do this
    before (c).
  \item \textbf{By measurement.} Build a program whose handler toggles a pin and does nothing
    else: make \code{PB0} an output, and let the handler be \code{sbi PINB, 0} and \code{reti}.
    Then:
\begin{shell}
make measure IMAGE=app CALLS="--watch B0 --run 200000"
\end{shell}
    \code{--watch} records every transition of a pin and reports the gaps between them, and a
    period is a gap. Ignore the first gap; part (h) is about why.
  \item \textbf{Reconcile.} All three should agree exactly, and they do. Say in one sentence why a
    timer is the one thing in this book where that is unsurprising.
  \item \textbf{Now ask for 3\,kHz instead.} Compute the prescaler, compare value and period in
    cycles, then measure it the same way.
  \item \textbf{The discrepancy.} Your computed period is 5333 cycles. \textbf{No measured gap is
    5333.} Write down what you actually measured, then answer three questions:
    \begin{itemize}
      \item Is the timer wrong?
      \item Is the average right?
      \item What is quantising the measurement, given that the timer is exact?
    \end{itemize}
  \item \textbf{Test your explanation.} Whatever you decided in (f), it makes a prediction about
    compare values that give an \textbf{even} period. Change \code{OCR1A} to 5331 and then to 5333,
    measure both, and say whether your explanation survived.
  \item \textbf{The first gap.} In every one of these runs the first gap is different from all the
    others: a few cycles longer than a period. Explain it, and say whether it would still happen on
    real hardware.
\end{parts}

% ------------------------------------------------------------------------------------------
\exercise{When the handler is too slow}{Design}
\label{c5:ex:slow}
\begin{parts}
  \item Your handler takes 800 cycles and the compare match happens every 16000. What fraction of
    the processor is the handler using?
  \item The same handler, with the interrupt rate raised to 100\,kHz. Now what?
  \item Describe precisely what happens when a compare match occurs while the previous handler is
    still running. Is the interrupt lost or delayed?
  \item Two matches occur during one long handler. Now what, and how does this differ from (c)?
  \item The software timers built on that interrupt now run slow. Say by how much, in terms of the
    interrupts that were missed, and say whether anything in the program could detect it.
  \item Give two ways to make the problem visible rather than silent, and say what each costs.
\end{parts}
